% Options for packages loaded elsewhere
% Options for packages loaded elsewhere
\PassOptionsToPackage{unicode}{hyperref}
\PassOptionsToPackage{hyphens}{url}
\PassOptionsToPackage{dvipsnames,svgnames,x11names}{xcolor}
%
\documentclass[
  letterpaper,
  DIV=11,
  numbers=noendperiod]{scrartcl}
\usepackage{xcolor}
\usepackage{amsmath,amssymb}
\setcounter{secnumdepth}{5}
\usepackage{iftex}
\ifPDFTeX
  \usepackage[T1]{fontenc}
  \usepackage[utf8]{inputenc}
  \usepackage{textcomp} % provide euro and other symbols
\else % if luatex or xetex
  \usepackage{unicode-math} % this also loads fontspec
  \defaultfontfeatures{Scale=MatchLowercase}
  \defaultfontfeatures[\rmfamily]{Ligatures=TeX,Scale=1}
\fi
\usepackage{lmodern}
\ifPDFTeX\else
  % xetex/luatex font selection
\fi
% Use upquote if available, for straight quotes in verbatim environments
\IfFileExists{upquote.sty}{\usepackage{upquote}}{}
\IfFileExists{microtype.sty}{% use microtype if available
  \usepackage[]{microtype}
  \UseMicrotypeSet[protrusion]{basicmath} % disable protrusion for tt fonts
}{}
\makeatletter
\@ifundefined{KOMAClassName}{% if non-KOMA class
  \IfFileExists{parskip.sty}{%
    \usepackage{parskip}
  }{% else
    \setlength{\parindent}{0pt}
    \setlength{\parskip}{6pt plus 2pt minus 1pt}}
}{% if KOMA class
  \KOMAoptions{parskip=half}}
\makeatother
% Make \paragraph and \subparagraph free-standing
\makeatletter
\ifx\paragraph\undefined\else
  \let\oldparagraph\paragraph
  \renewcommand{\paragraph}{
    \@ifstar
      \xxxParagraphStar
      \xxxParagraphNoStar
  }
  \newcommand{\xxxParagraphStar}[1]{\oldparagraph*{#1}\mbox{}}
  \newcommand{\xxxParagraphNoStar}[1]{\oldparagraph{#1}\mbox{}}
\fi
\ifx\subparagraph\undefined\else
  \let\oldsubparagraph\subparagraph
  \renewcommand{\subparagraph}{
    \@ifstar
      \xxxSubParagraphStar
      \xxxSubParagraphNoStar
  }
  \newcommand{\xxxSubParagraphStar}[1]{\oldsubparagraph*{#1}\mbox{}}
  \newcommand{\xxxSubParagraphNoStar}[1]{\oldsubparagraph{#1}\mbox{}}
\fi
\makeatother


\usepackage{longtable,booktabs,array}
\newcounter{none} % for unnumbered tables
\usepackage{calc} % for calculating minipage widths
% Correct order of tables after \paragraph or \subparagraph
\usepackage{etoolbox}
\makeatletter
\patchcmd\longtable{\par}{\if@noskipsec\mbox{}\fi\par}{}{}
\makeatother
% Allow footnotes in longtable head/foot
\IfFileExists{footnotehyper.sty}{\usepackage{footnotehyper}}{\usepackage{footnote}}
\makesavenoteenv{longtable}
\usepackage{graphicx}
\makeatletter
\newsavebox\pandoc@box
\newcommand*\pandocbounded[1]{% scales image to fit in text height/width
  \sbox\pandoc@box{#1}%
  \Gscale@div\@tempa{\textheight}{\dimexpr\ht\pandoc@box+\dp\pandoc@box\relax}%
  \Gscale@div\@tempb{\linewidth}{\wd\pandoc@box}%
  \ifdim\@tempb\p@<\@tempa\p@\let\@tempa\@tempb\fi% select the smaller of both
  \ifdim\@tempa\p@<\p@\scalebox{\@tempa}{\usebox\pandoc@box}%
  \else\usebox{\pandoc@box}%
  \fi%
}
% Set default figure placement to htbp
\def\fps@figure{htbp}
\makeatother





\setlength{\emergencystretch}{3em} % prevent overfull lines

\providecommand{\tightlist}{%
  \setlength{\itemsep}{0pt}\setlength{\parskip}{0pt}}



 


\KOMAoption{captions}{tableheading}
\makeatletter
\@ifpackageloaded{caption}{}{\usepackage{caption}}
\AtBeginDocument{%
\ifdefined\contentsname
  \renewcommand*\contentsname{Table of contents}
\else
  \newcommand\contentsname{Table of contents}
\fi
\ifdefined\listfigurename
  \renewcommand*\listfigurename{List of Figures}
\else
  \newcommand\listfigurename{List of Figures}
\fi
\ifdefined\listtablename
  \renewcommand*\listtablename{List of Tables}
\else
  \newcommand\listtablename{List of Tables}
\fi
\ifdefined\figurename
  \renewcommand*\figurename{Figure}
\else
  \newcommand\figurename{Figure}
\fi
\ifdefined\tablename
  \renewcommand*\tablename{Table}
\else
  \newcommand\tablename{Table}
\fi
}
\@ifpackageloaded{float}{}{\usepackage{float}}
\floatstyle{ruled}
\@ifundefined{c@chapter}{\newfloat{codelisting}{h}{lop}}{\newfloat{codelisting}{h}{lop}[chapter]}
\floatname{codelisting}{Listing}
\newcommand*\listoflistings{\listof{codelisting}{List of Listings}}
\makeatother
\makeatletter
\makeatother
\makeatletter
\@ifpackageloaded{caption}{}{\usepackage{caption}}
\@ifpackageloaded{subcaption}{}{\usepackage{subcaption}}
\makeatother
\usepackage{bookmark}
\IfFileExists{xurl.sty}{\usepackage{xurl}}{} % add URL line breaks if available
\urlstyle{same}
\hypersetup{
  colorlinks=true,
  linkcolor={blue},
  filecolor={Maroon},
  citecolor={Blue},
  urlcolor={Blue},
  pdfcreator={LaTeX via pandoc}}


\author{}
\date{}
\begin{document}


\section{Tabla de Resumen}\label{tabla-de-resumen}

\begin{landscape}

\begin{longtable}[]{@{}
  >{\raggedright\arraybackslash}p{(\linewidth - 10\tabcolsep) * \real{0.1667}}
  >{\raggedright\arraybackslash}p{(\linewidth - 10\tabcolsep) * \real{0.1667}}
  >{\raggedright\arraybackslash}p{(\linewidth - 10\tabcolsep) * \real{0.1667}}
  >{\centering\arraybackslash}p{(\linewidth - 10\tabcolsep) * \real{0.1667}}
  >{\raggedright\arraybackslash}p{(\linewidth - 10\tabcolsep) * \real{0.1667}}
  >{\raggedright\arraybackslash}p{(\linewidth - 10\tabcolsep) * \real{0.1667}}@{}}
\caption{Tabla 1. Agrupación de todos los indicadores de desarrollo
humano, pobreza y desigualdad.}\label{tbl-resumen}\tabularnewline
\toprule\noalign{}
\begin{minipage}[b]{\linewidth}\raggedright
Indicador
\end{minipage} & \begin{minipage}[b]{\linewidth}\raggedright
Valor Provincial
\end{minipage} & \begin{minipage}[b]{\linewidth}\raggedright
Referente Nacional
\end{minipage} & \begin{minipage}[b]{\linewidth}\centering
Año
\end{minipage} & \begin{minipage}[b]{\linewidth}\raggedright
Fuente
\end{minipage} & \begin{minipage}[b]{\linewidth}\raggedright
Interpretación Breve
\end{minipage} \\
\midrule\noalign{}
\endfirsthead
\toprule\noalign{}
\begin{minipage}[b]{\linewidth}\raggedright
Indicador
\end{minipage} & \begin{minipage}[b]{\linewidth}\raggedright
Valor Provincial
\end{minipage} & \begin{minipage}[b]{\linewidth}\raggedright
Referente Nacional
\end{minipage} & \begin{minipage}[b]{\linewidth}\centering
Año
\end{minipage} & \begin{minipage}[b]{\linewidth}\raggedright
Fuente
\end{minipage} & \begin{minipage}[b]{\linewidth}\raggedright
Interpretación Breve
\end{minipage} \\
\midrule\noalign{}
\endhead
\bottomrule\noalign{}
\endlastfoot
Índice de Desarrollo Humano (IDH) & 0,78 & 0,77 & 2025 & \textbf{Val.
Prov.:} Datos ENEMDU \textbf{Ref. Nal.:} PNUD & Podemos observar que
Pichincha presenta un desarrollo humano alto, con un valor ligeramente
superior al promedio nacional, donde sus mayores fortalezas corresponden
a las dimensiones de salud y educación, mientras que el ingreso
constituye la principal limitación para alcanzar un mayor nivel de
desarrollo. \\
IDH - Modificado & 0,81 & No existe & 2025 & \textbf{Val. Prov.:} ENEMDU
2025 e INEC. \textbf{Ref. Nal.:} No existe & Al incorporar la dimensión
de migración, el índice aumenta respecto al IDH tradicional, lo que
evidencia que la dinámica migratoria aporta favorablemente al desarrollo
humano de la provincia. Al tratarse de un indicador construido para esta
investigación, no existe un referente nacional oficial para su
comparación. \\
Índice de Pobreza Multidimensional (IPM) & & & 2025 & \textbf{Val.
Prov.:} Datos ENEMDU \textbf{Ref. Nal.:} & \\
IPM - Modificado & & & 2025 & \textbf{Val. Prov.:} \textbf{Ref. Nal.:}
& \\
Pobreza por ingresos & & & 2025 & \textbf{Val. Prov.:} \textbf{Ref.
Nal.:} & \\
Índice de Theil & 0,39 & 0,39 & 2025 & \textbf{Val. Prov.:} Datos ENEMDU
\textbf{Ref. Nal.:} Laboratorio de Equidad del Banco Mundial & El valor
tanto a nivel nacional como provincial es el mismo y este se encuentra
en una escala de desigualdad alta moderada, lo que significa que no hay
una distribución equitativa del ingreso en el territorio, es decir, que
una parte de la población concentra ingresos superiores al promedio,
mientras otros grupos se mantienen rezagados. \\
Curva de Lorenz & Ver Figura 1. & Ver Figura 2. & 2025 & \textbf{Val.
Prov.:} Datos ENEMDU \textbf{Ref. Nal.:} Narriaga, K. (oct. 2025) & A
pesar de que ambas curvas se mantienen por debajo de la línea de
pobreza, la curva a nivel provincial tiene porcentajes de acumulación de
ingresos menores, y en ambos casos el 60\% de la población acumula
apenas el 30\% de los ingresos, acentuando aún más la desigualdad
económica que existe. \\
Coeficiente de Gini & 0,45 & 0,44 & 2025 & \textbf{Val. Prov.:} Datos
ENEMDU \textbf{Ref. Nal.:} INEC & Aunque se puede ver una pequeña
diferencia porcentual entre el coeficiente provincial y nacional, indica
que la desigualdad provincial es ligeramente mayor que la del país,
evidenciando que Pichincha concentra desigualdades en la distribución
del ingreso. \\
IDH Ajustado por Desigualdad (IDH-D) & 0,60 & & 2025 & \textbf{Val.
Prov.:} ENEMDU 2025, metodología PNUD. \textbf{Ref. Nal.:} & \\
\end{longtable}

\end{landscape}




\end{document}
